\documentclass[12pt,oneside]{book}

% ============================================================
% METADATOS
% ============================================================
\newcommand{\universidad}{Universidad de Buenos Aires}
\newcommand{\facultad}{Facultad de Derecho}
\newcommand{\materia}{Derecho Privado VI --- Derechos Reales}
\newcommand{\materiacorta}{Derecho Privado VI}
\newcommand{\titulo}{Garantías Reales: Hipoteca, Prenda y Anticresis}
\newcommand{\autor}{Isaac Edgar Camacho Ocampo}
\newcommand{\anio}{2025}

% ============================================================
% PAQUETES
% ============================================================
\usepackage[utf8]{inputenc}
\usepackage[T1]{fontenc}
\usepackage[spanish,es-nodecimaldot]{babel}

\usepackage[
    top=2.5cm,
    bottom=2.5cm,
    left=3cm,
    right=2.5cm,
    headheight=15pt
]{geometry}

\usepackage{amsmath}
\usepackage{amssymb}
\usepackage{mathtools}

\usepackage{graphicx}
\usepackage{float}
\usepackage{caption}
\usepackage{subcaption}

\usepackage{booktabs}
\usepackage{array}
\usepackage{longtable}
\usepackage{tabularx}

\usepackage{enumitem}

\usepackage{xcolor}
\usepackage[most]{tcolorbox}

\usepackage{fancyhdr}
\usepackage{ragged2e}

\usepackage[
    colorlinks=true,
    linkcolor=blue!50!black,
    citecolor=green!40!black,
    urlcolor=blue!60!black,
    pdftitle={Apuntes de Derecho Privado VI - Derechos Reales: Garantias Reales: Hipoteca, Prenda y Anticresis},
    pdfauthor={Isaac Edgar Camacho Ocampo},
    pdfsubject={Apuntes universitarios},
    bookmarks=true
]{hyperref}

% ============================================================
% CONFIGURACIÓN
% ============================================================
\graphicspath{
    {./img/}
    {./graficos/}
    {./figuras/}
}

\setcounter{tocdepth}{2}

\setlength{\parindent}{0pt}
\setlength{\parskip}{0.5em}

\setlist[itemize]{
    topsep=0.3em,
    itemsep=0.2em
}

\setlist[enumerate]{
    topsep=0.3em,
    itemsep=0.2em
}

\clubpenalty=10000
\widowpenalty=10000

\pagestyle{fancy}
\fancyhf{}
\fancyhead[L]{\nouppercase{\leftmark}}
\fancyhead[R]{\materiacorta}
\fancyfoot[C]{\thepage}
\renewcommand{\headrulewidth}{0.4pt}
\renewcommand{\footrulewidth}{0pt}

\fancypagestyle{plain}{
    \fancyhf{}
    \fancyfoot[C]{\thepage}
    \renewcommand{\headrulewidth}{0pt}
}

% ============================================================
% COMANDOS
% ============================================================
\newcommand{\abs}[1]{\left|#1\right|}
\newcommand{\norm}[1]{\left\lVert#1\right\rVert}
\newcommand{\vect}[1]{\mathbf{#1}}
\newcommand{\deriv}[2]{\frac{d#1}{d#2}}
\newcommand{\pderiv}[2]{\frac{\partial#1}{\partial#2}}

% ============================================================
% ENTORNOS / CAJAS
% ============================================================
\newtcolorbox{definition}[1]{
    enhanced,
    breakable,
    title={Definición: #1},
    fonttitle=\bfseries,
    colback=blue!3,
    colframe=blue!50!black,
    boxrule=0.8pt,
    arc=2mm
}

\newtcolorbox{important}{
    enhanced,
    breakable,
    title={Importante},
    fonttitle=\bfseries,
    colback=yellow!8,
    colframe=orange!70!black,
    boxrule=0.8pt,
    arc=2mm
}

\newtcolorbox{example}[1]{
    enhanced,
    breakable,
    title={Ejemplo: #1},
    fonttitle=\bfseries,
    colback=green!4,
    colframe=green!40!black,
    boxrule=0.8pt,
    arc=2mm
}

\newtcolorbox{formula}[1]{
    enhanced,
    breakable,
    title={Fórmula / Regla: #1},
    fonttitle=\bfseries,
    colback=purple!4,
    colframe=purple!50!black,
    boxrule=0.8pt,
    arc=2mm
}

\newtcolorbox{exercise}[1]{
    enhanced,
    breakable,
    title={Ejercicio: #1},
    fonttitle=\bfseries,
    colback=gray!5,
    colframe=gray!60!black,
    boxrule=0.8pt,
    arc=2mm
}

\newtcolorbox{note}{
    enhanced,
    breakable,
    title={Nota},
    fonttitle=\bfseries,
    colback=white,
    colframe=gray!50,
    boxrule=0.6pt,
    arc=2mm
}

% ============================================================
% CONFIGURACIÓN FINAL
% ============================================================
\renewcommand{\arraystretch}{1.25}

% ============================================================
\begin{document}
% ============================================================

% ------------------------------------------------------------
% PORTADA
% ------------------------------------------------------------
\begin{titlepage}

\thispagestyle{empty}

\begin{center}

\vspace*{1cm}

{\large \textsc{\universidad}}

\vspace{0.2cm}

{\large \textsc{\facultad}}

\vspace{3cm}

{\Huge \textbf{\materia}}

\vspace{1cm}

{\LARGE \titulo}

\vspace{0.8cm}

\textit{Estudiar con constancia es convertir cada concepto en una herramienta propia.}

\vspace{1.2cm}

\rule{0.70\textwidth}{0.5pt}

\vspace{0.5cm}

{\large Apuntes de estudio}

\vspace{0.5cm}

\rule{0.70\textwidth}{0.5pt}

\vfill

{\large \autor}

\vspace{0.3cm}

{\large \anio}

\end{center}

\vfill

\begin{flushleft}
\textit{Este es un modesto aporte a los estudiantes de ``Derecho Privado VI --- Derechos Reales'' no pretende sustituir el material oficial de la materia.}
\end{flushleft}

\end{titlepage}

% ------------------------------------------------------------
% ÍNDICE
% ------------------------------------------------------------
\tableofcontents

% ------------------------------------------------------------
% CONTENIDO
% ------------------------------------------------------------

% ============================================================
\chapter[Marco general de las garantías reales]{Marco general de las garantías reales}
% ============================================================

\section{Presentación del tema}

En esta clase se abordan los tres derechos reales de garantía regulados por el Código Civil y Comercial de la Nación (Ley 26.994): la \textbf{hipoteca}, la \textbf{prenda} y la \textbf{anticresis}. Se analizan sus diferencias esenciales en cuanto al objeto, la posesión y los legitimados para constituirlos, partiendo de las disposiciones comunes ya estudiadas y profundizando en la normativa específica de cada figura.

Las garantías reales aparecen en la práctica profesional del abogado en operaciones que involucran hipotecas (compraventa de inmuebles financiadas, créditos bancarios), prendas (adquisición de automotores, maquinaria industrial) y anticresis (contratos de locación con garantía real). Su comprensión es necesaria para asesorar en:

\begin{itemize}
    \item Operaciones inmobiliarias: compraventa, hipoteca, financiamiento bancario.
    \item Contratos de crédito con garantía real.
    \item Ejecuciones hipotecarias y prendarias (subasta judicial).
    \item Cancelación registral de gravámenes y liberación de inmuebles.
    \item Asesoramiento a condóminos, superficiarios y usufructuarios.
    \item Redacción de contratos de locación con garantía de prenda de moneda extranjera.
\end{itemize}

\begin{example}{Relación entre los bloques}
Se plantea el caso de quien desea comprar un departamento, un auto y, además, tiene un campo productivo que quiere usar como garantía de un crédito. Según el bien y la forma en que se quiera estructurar la garantía, el código ofrece tres figuras distintas: hipoteca para el departamento, prenda para el auto y anticresis para el campo. Todas comparten los mismos fundamentos normativos, pero cada una tiene sus propias reglas según el objeto y la posesión.
% TODO: contenido incompleto o ambiguo en las notas originales (el apunte asigna prenda al auto y anticresis al campo, aunque más adelante señala que el automotor se prenda mediante prenda con registro)
\end{example}

\section{Disposiciones comunes a todas las garantías reales}

En las clases anteriores se trataron las disposiciones comunes en materia de garantías reales. En esta clase se aborda el tratamiento de sus diferencias y de las características propias de cada uno de estos derechos accesorios de garantía.

Los derechos reales de garantía son tres: \textbf{hipoteca}, \textbf{prenda} y \textbf{anticresis}. A los tres se les aplica necesariamente toda la normativa general en materia de derechos reales.

\begin{important}
\textbf{Normativa general aplicable a todos los derechos reales:} arts.\ 1882 a 1907 del Código Civil y Comercial de la Nación.

También rigen las normas propias aplicables a todas las garantías reales (\textit{disposiciones comunes en materia de garantías}): arts.\ 2184 a 2204.
\end{important}

A partir de los artículos mencionados, se aplican las normas propias de cada garantía:

\begin{itemize}
    \item \textbf{Hipoteca:} arts.\ 2205 a 2211.
    \item \textbf{Anticresis:} arts.\ 2212 a 2218.
    \item \textbf{Prenda:} existe una novedad, pues el codificador dedica artículos a disposiciones comunes aplicables a los dos tipos de prenda (prenda de cosas y prenda de créditos instrumentados), antes de las normas específicas.
\end{itemize}

\section{Caracteres de las garantías reales}

Las disposiciones comunes en materia de garantías (arts.\ 2184 a 2204) establecen los siguientes caracteres esenciales:

\begin{itemize}
    \item Convencionalidad.
    \item Accesoriedad.
    \item Especialidad objetiva.
    \item Especialidad crediticia.
    \item Indivisibilidad.
\end{itemize}

\begin{important}
Estos caracteres deben estar presentes necesariamente en todas las garantías reales, \textbf{bajo pena de nulidad}.
\end{important}

\section{Diferencias esenciales: el objeto de la garantía}

Las tres garantías comparten muchas disposiciones y se les aplican las mismas normas. Sin embargo, se diferencian esencialmente en cuanto al \textbf{objeto} sobre el que recae el derecho real, es decir, el objeto que soporta el privilegio del acreedor y que será ejecutado si el deudor no cumple con la obligación. Esta diferencia permite identificar con claridad el tipo de garantía.

\begin{table}[H]
\centering
\begin{tabularx}{\textwidth}{
    >{\raggedright\arraybackslash}p{0.22\textwidth}
    >{\raggedright\arraybackslash}X
}
\toprule
\textbf{Garantía} & \textbf{Objeto} \\
\midrule
Hipoteca & Un inmueble. También en el caso del derecho de superficie. \\
Prenda & Cosas muebles no registrables o créditos instrumentados. \\
Anticresis & Cosas registrables (pueden ser muebles o inmuebles). \\
\bottomrule
\end{tabularx}
\end{table}

\section{La posesión como criterio diferenciador}

La posesión constituye otra de las grandes diferencias entre las garantías reales.

\begin{table}[H]
\centering
\begin{tabularx}{\textwidth}{
    >{\raggedright\arraybackslash}p{0.20\textwidth}
    >{\raggedright\arraybackslash}X
    >{\raggedright\arraybackslash}p{0.27\textwidth}
}
\toprule
\textbf{Garantía} & \textbf{¿Quién tiene la posesión?} & \textbf{Publicidad} \\
\midrule
Hipoteca & El constituyente (titular del dominio). El acreedor \textbf{no} tiene posesión. & Registro de la Propiedad Inmueble. \\
Prenda clásica & Se desplaza al acreedor (quien la recupera al extinguirse la obligación). & La posesión misma. \\
Anticresis & Se desplaza al acreedor o al tercero designado por las partes. & La posesión de la cosa. \\
\bottomrule
\end{tabularx}
\end{table}

\begin{example}{Posesión en la hipoteca}
Si una persona es propietaria de un inmueble y tramitó un crédito con garantía hipotecaria (por ejemplo, por saldo de precio en el Banco Nación), no es el banco quien tiene la posesión del departamento. El propietario está en posesión del inmueble y es el titular dominial de la cosa hipotecada.
\end{example}

\subsection*{Cuadro Cornell: marco general}

\begin{longtable}{p{0.24\textwidth}p{0.69\textwidth}}
\toprule
\textbf{Preguntas / palabras clave} & \textbf{Notas / respuestas} \\
\midrule
\endfirsthead
\toprule
\textbf{Preguntas / palabras clave} & \textbf{Notas / respuestas} \\
\midrule
\endhead

\textbf{¿Cuáles son los tres derechos reales de garantía?} &
Hipoteca, prenda y anticresis. A los tres les son aplicables las normas generales de derechos reales (arts.\ 1882 a 1907 CCCN) y las disposiciones comunes en materia de garantías (arts.\ 2184 a 2204 CCCN). \\[0.6em]

\textbf{¿Qué normas específicas regulan cada garantía?} &
Hipoteca: arts.\ 2205 a 2211. Anticresis: arts.\ 2212 a 2218. Prenda: disposiciones comunes a los dos tipos de prenda (prenda de cosas y de créditos instrumentados), seguidas de las normas específicas. \\[0.6em]

\textbf{¿Qué caracteres deben estar presentes en toda garantía real?} &
Convencionalidad, accesoriedad, especialidad objetiva, especialidad crediticia e indivisibilidad, bajo pena de nulidad. \\[0.6em]

\textbf{¿En qué se diferencian las tres garantías en cuanto al objeto?} &
La hipoteca recae sobre un inmueble (también el derecho de superficie); la prenda, sobre cosas muebles no registrables o créditos instrumentados; la anticresis, sobre cosas registrables, muebles o inmuebles. \\[0.6em]

\textbf{¿Quién tiene la posesión en cada garantía y cómo se da la publicidad?} &
En la hipoteca la conserva el constituyente y la publicidad es registral (Registro de la Propiedad Inmueble). En la prenda clásica se desplaza al acreedor y la publicidad es la posesión. En la anticresis se desplaza al acreedor o al tercero designado y la publicidad es la posesión de la cosa. \\

\midrule
\multicolumn{2}{p{\dimexpr0.93\textwidth+2\tabcolsep\relax}}{%
\textbf{Resumen:} los tres derechos reales de garantía comparten un mismo sustrato normativo (normas generales de los derechos reales y disposiciones comunes de las garantías), pero se diferencian esencialmente en el objeto sobre el que recaen y en el régimen de posesión.} \\
\bottomrule
\end{longtable}

% ============================================================
\chapter[Hipoteca]{Hipoteca}
% ============================================================

\section{Concepto de hipoteca}

\begin{definition}{Hipoteca (art.\ 2205 CCCN)}
``La hipoteca es un derecho real que recae sobre uno o más inmuebles individualizados que continúan en poder del constituyente y que otorgan al acreedor, ante el incumplimiento del deudor, las facultades de persecución y preferencia para cobrar con su producido el crédito garantizado.''
\end{definition}

El artículo 2205 es claro en que el derecho real de garantía recae sobre uno o más inmuebles individualizados que \textbf{continúan en poder del constituyente}. La posesión no se desplaza al acreedor; el constituyente es siempre el propietario del inmueble. Ese propietario en ocasiones es deudor y en ocasiones no lo es: tendrá la calidad de \textbf{propietario no deudor}.

La fortaleza de este derecho real es la \textbf{persecución} y la \textbf{preferencia}. Esto permite, por ejemplo, que el titular de dominio transmita el inmueble sin que ello afecte al acreedor, quien persigue la cosa en manos de quien se encuentre, aun cuando hubiese cambiado el titular dominial.

La persecución y la preferencia se ejercen ante el incumplimiento del deudor. En ese momento, el acreedor intima el pago:

\begin{itemize}
    \item Si el deudor paga, la hipoteca se extingue (aunque el inmueble continúe figurando como hipotecado en el registro hasta la cancelación).
    \item Si el deudor no paga, se pone en marcha el procedimiento de ejecución: la subasta del inmueble.
\end{itemize}

\section{Legitimados para constituir hipoteca}

\begin{definition}{Legitimados (art.\ 2206 CCCN)}
Únicamente pueden constituir hipoteca:
\begin{itemize}
    \item El \textbf{titular del dominio} (dueño pleno), que tiene el \textit{ius disponendi} que le permite enajenar y constituir derechos reales de disfrute o de garantía.
    \item Los \textbf{condóminos}, respecto de su parte indivisa.
    \item Quienes sean propietarios en \textbf{propiedad horizontal}.
    \item Propietarios en \textbf{conjuntos inmobiliarios}.
    \item El \textbf{superficiario} (quien tiene derecho de construir, plantar o forestar sobre inmueble ajeno), que puede hipotecar su derecho.
\end{itemize}
\end{definition}

\subsection{Hipoteca de la parte indivisa en condominio}

El condómino no puede hipotecar toda la cosa, porque sus facultades respecto de ella son restringidas. Sí puede hipotecar su \textbf{parte indivisa}. Esa parte indivisa puede ser ejecutada (subastada) por el acreedor en caso de incumplimiento. El adquirente en subasta adquiriría, por ejemplo, un 33,33\,\% del inmueble y pasaría a ser condómino junto con los otros titulares.

\begin{important}
Mientras la hipoteca sobre la parte indivisa subsista, si los condóminos realizan la partición extrajudicial, el acreedor necesariamente deberá ser parte de ese acto. De lo contrario, la partición le será \textbf{inoponible} al acreedor.
\end{important}

\section{Forma del contrato constitutivo}

Se trata de un inmueble. Como en materia de derechos reales sobre inmuebles todo debe realizarse necesariamente en escritura pública, así ocurre también en la hipoteca.

\begin{important}
La hipoteca se constituye por \textbf{escritura pública}. Del mismo modo, para cancelar la inscripción registral también se requiere escritura pública. Cuando esa escritura pública de cancelación se inscribe en el Registro de la Propiedad Inmueble, opera la cancelación.
\end{important}

\section{Efectos del registro: plazo de inscripción}

\begin{important}
La Ley 27.271 sustituyó los artículos 2189 y 2210 del Código Civil y Comercial. Esta reforma no se encuentra en el texto de la mayoría de las ediciones del código que se utilizan en la práctica, por lo que es fundamental tenerla en cuenta.
\end{important}

\begin{formula}{Duración de la inscripción hipotecaria (art.\ 2210, según Ley 27.271)}
Los efectos del registro de la hipoteca se conservan por el término de \textbf{35 años} si antes no se renueva.
\end{formula}

Evolución normativa del plazo de inscripción:

\begin{table}[H]
\centering
\begin{tabularx}{\textwidth}{
    >{\raggedright\arraybackslash}p{0.34\textwidth}
    >{\raggedright\arraybackslash}p{0.18\textwidth}
    >{\raggedright\arraybackslash}X
}
\toprule
\textbf{Período} & \textbf{Plazo} & \textbf{Fuente normativa} \\
\midrule
Código 2015 hasta el 15/09/2016 & 20 años & Texto original del CCCN (art.\ 2210) \\
Desde el 15 de septiembre de 2016 en adelante & 35 años & Ley 27.271 (reforma del art.\ 2210) \\
\bottomrule
\end{tabularx}
\end{table}

\section{Extinción y cancelación de la hipoteca}

Es fundamental distinguir estas dos situaciones.

\begin{table}[H]
\centering
\begin{tabularx}{\textwidth}{
    >{\raggedright\arraybackslash}p{0.17\textwidth}
    >{\raggedright\arraybackslash}p{0.30\textwidth}
    >{\raggedright\arraybackslash}X
}
\toprule
\textbf{Concepto} & \textbf{¿Cuándo ocurre?} & \textbf{Efecto} \\
\midrule
\textbf{Extinción} & Cuando se extingue la obligación principal (por ejemplo, pago de la última cuota del crédito). & La hipoteca se extingue por vía de consecuencia (carácter accesorio). Pero el inmueble puede seguir figurando como hipotecado en el registro. \\
\textbf{Cancelación} & Cuando se otorga escritura pública de cancelación y se la inscribe en el Registro de la Propiedad Inmueble. & Recién aquí el gravamen desaparece frente a terceros. Sin cancelación, el inmueble sigue hipotecado ante el registro. \\
\bottomrule
\end{tabularx}
\end{table}

\begin{example}{Hipotecas extinguidas pero no canceladas}
En la práctica profesional hay numerosos casos en que las hipotecas se encuentran extinguidas pero no canceladas, y el propietario lo desconoce. Cuando va a vender el inmueble, allí aparece la hipoteca. Por eso es tan importante distinguir extinción de cancelación.
\end{example}

\begin{note}
\textbf{Seguros de vida en créditos bancarios.} En los créditos bancarios en general existe un seguro para extinguir las obligaciones ante el fallecimiento del deudor. Pero este seguro no comprende la cancelación registral. Por lo tanto, el seguro extingue la obligación, pero ante terceros queda pendiente la cancelación hasta que los herederos realicen el trámite necesario para otorgar la escritura y cancelar la inscripción.
\end{note}

\subsection{Tipos de cancelación}

\begin{itemize}
    \item \textbf{Cancelación voluntaria:} cuando el propio acreedor (particular o bancario) realiza la escritura de cancelación.
    \item \textbf{Cancelación judicial:} cuando el acreedor no cancela voluntariamente. Requiere un juicio específico de cancelación de hipoteca, en el que el constituyente demanda al acreedor.
    \item \textbf{Cancelación de pleno derecho:} opera por mero efecto de la ley, cuando vence el plazo de inscripción sin que se haya producido la renovación. Antes: 20 años; actualmente: 35 años (según la reforma de la Ley 27.271).
\end{itemize}

\section{Ejecución hipotecaria}

La ejecución de la hipoteca (subasta del inmueble ante el incumplimiento del deudor) está regulada en los Códigos de Procedimiento de cada jurisdicción, no directamente en el Código Civil y Comercial.

\subsection*{Cuadro Cornell: hipoteca}

\begin{longtable}{p{0.24\textwidth}p{0.69\textwidth}}
\toprule
\textbf{Preguntas / palabras clave} & \textbf{Notas / respuestas} \\
\midrule
\endfirsthead
\toprule
\textbf{Preguntas / palabras clave} & \textbf{Notas / respuestas} \\
\midrule
\endhead

\textbf{¿Qué es la hipoteca y sobre qué recae?} &
Es el derecho real que recae sobre uno o más inmuebles individualizados que continúan en poder del constituyente y que otorga al acreedor, ante el incumplimiento del deudor, las facultades de persecución y preferencia para cobrar con su producido el crédito garantizado (art.\ 2205). El derecho recae sobre inmuebles (o el derecho de superficie). \\[0.6em]

\textbf{¿Qué significa que la posesión no se desplace al acreedor?} &
El constituyente, que siempre es el propietario del inmueble, conserva la posesión; el acreedor no la tiene. La publicidad se otorga mediante la inscripción en el Registro de la Propiedad Inmueble. El propietario puede ser deudor o propietario no deudor. \\[0.6em]

\textbf{¿Qué permiten la persecución y la preferencia?} &
Que el titular de dominio transmita el inmueble sin afectar al acreedor, quien persigue la cosa en manos de quien se encuentre. Se ejercen ante el incumplimiento: el acreedor intima el pago; si el deudor paga, la hipoteca se extingue; si no paga, se subasta el inmueble. \\[0.6em]

\textbf{¿Quiénes pueden constituir hipoteca?} &
El titular del dominio pleno, los condóminos (solo sobre su parte indivisa), los propietarios en propiedad horizontal o en conjuntos inmobiliarios, y el superficiario respecto de su derecho (art.\ 2206). \\[0.6em]

\textbf{¿Qué ocurre con la parte indivisa hipotecada en un condominio?} &
El condómino no puede hipotecar toda la cosa, pero sí su parte indivisa, que puede ser subastada; el adquirente pasa a ser condómino. Si hay partición extrajudicial mientras la hipoteca subsiste, el acreedor debe ser parte del acto; de lo contrario la partición le es inoponible. \\[0.6em]

\textbf{¿Qué forma exige la hipoteca?} &
Escritura pública, tanto para constituirla como para cancelar la inscripción registral; la cancelación opera cuando la escritura de cancelación se inscribe en el Registro de la Propiedad Inmueble. \\[0.6em]

\textbf{¿Cuál es el plazo de inscripción hipotecaria vigente?} &
35 años, desde la reforma de la Ley 27.271 (que sustituyó los arts.\ 2189 y 2210), a partir del 15/09/2016. Antes, el plazo era de 20 años. \\[0.6em]

\textbf{¿Qué diferencia hay entre extinción y cancelación?} &
La extinción ocurre al extinguirse la obligación principal (p.\,ej., pago de la última cuota), por el carácter accesorio. La cancelación requiere escritura pública inscripta en el registro; sin ella, el inmueble sigue hipotecado frente a terceros aunque la deuda esté pagada. El seguro de vida de los créditos bancarios extingue la obligación pero no cancela el registro. \\[0.6em]

\textbf{¿Cuáles son las formas de cancelación?} &
Voluntaria (escritura realizada por el acreedor), judicial (juicio de cancelación en que el constituyente demanda al acreedor) y de pleno derecho (vencimiento del plazo de inscripción sin renovación). \\[0.6em]

\textbf{¿Dónde se regula la ejecución hipotecaria?} &
En los Códigos de Procedimiento de cada jurisdicción, no directamente en el Código Civil y Comercial. \\

\midrule
\multicolumn{2}{p{\dimexpr0.93\textwidth+2\tabcolsep\relax}}{%
\textbf{Resumen:} la hipoteca opera sobre inmuebles sin desplazamiento posesorio y otorga publicidad a través del Registro de la Propiedad Inmueble. Se constituye por escritura pública, sus efectos registrales se conservan por 35 años (Ley 27.271) y debe distinguirse rigurosamente la extinción de la cancelación, ya que la primera no implica automáticamente la segunda.} \\
\bottomrule
\end{longtable}

% ============================================================
\chapter[Anticresis]{Anticresis}
% ============================================================

\section{Concepto y objeto}

El derecho real de anticresis recae sobre cosas registrables individualizadas, sean muebles o inmuebles. A diferencia de la hipoteca, en la anticresis la posesión se entrega al acreedor (o a un tercero designado por las partes), ya que es esa posesión la que otorga publicidad, en lugar del registro.

\begin{definition}{Anticresis (arts.\ 2212 a 2218 CCCN)}
La anticresis es el derecho real de garantía que recae sobre cosas registrables (muebles o inmuebles). El acreedor recibe la posesión de la cosa como garantía. La publicidad la confiere la posesión, a diferencia de la hipoteca, donde la publicidad la confiere el registro.
\end{definition}

\section{Legitimados para constituir anticresis}

\begin{definition}{Legitimados (art.\ 2213 CCCN)}
Pueden constituir anticresis:
\begin{itemize}
    \item El titular del dominio.
    \item Los condóminos.
    \item Los propietarios en propiedad horizontal.
    \item El superficiario.
    \item El \textbf{usufructuario}: puede constituir anticresis, ya que tiene el derecho de disponer de los frutos, pero con la limitación temporal de su propio derecho (si es vitalicio, hasta su muerte; si es a término, hasta que se cumpla el plazo de 10 o 20 años según corresponda).
\end{itemize}
\end{definition}

\section{Plazos máximos}

La anticresis tiene un plazo máximo contado desde la constitución.

\begin{table}[H]
\centering
\begin{tabularx}{\textwidth}{
    >{\raggedright\arraybackslash}p{0.35\textwidth}
    >{\raggedright\arraybackslash}X
}
\toprule
\textbf{Tipo de bien} & \textbf{Plazo máximo} \\
\midrule
Muebles registrables & 5 años desde la constitución. \\
Inmuebles & 10 años desde la constitución. \\
\bottomrule
\end{tabularx}
\end{table}

\begin{note}
A diferencia de lo que ocurre con la hipoteca, no se ha prorrogado el plazo de la anticresis de 20 a 35 años (conf.\ art.\ 2218). El artículo 2218 mantiene las duraciones de 5 y 10 años según el tipo de bien.
\end{note}

\section{Derechos y deberes del acreedor anticresista}

El acreedor anticresista, al tener la posesión del objeto de la garantía, adquiere los siguientes derechos:

\begin{itemize}
    \item \textbf{Usar la cosa:} puede usar y explotar el bien, inclusive darlo en arrendamiento.
    \item \textbf{Percibir los frutos:} si es un inmueble, los frutos son civiles (rentas, alquileres). Se imputan primero a gastos e intereses y luego a capital.
\end{itemize}

Tiene, además, los siguientes deberes fundamentales:

\begin{itemize}
    \item \textbf{Conservación de la cosa:} está obligado a conservar y mantener el bien.
    \item \textbf{Pago de gastos de conservación:} debe afrontar los gastos de conservación de la cosa.
    \item \textbf{Recupero de mejoras:} puede recuperar el costo de las mejoras necesarias. Las mejoras útiles, solo en los casos en que le hubiesen conferido a la cosa un mayor valor.
    \item \textbf{Restitución:} al cumplimiento del plazo, o antes si se extingue la obligación, debe restituir la cosa al constituyente.
\end{itemize}

Como se trata de cosas registrables, en todos los casos se procede a la inscripción en el registro correspondiente (ya sea de muebles o de inmuebles).

\subsection*{Cuadro Cornell: anticresis}

\begin{longtable}{p{0.24\textwidth}p{0.69\textwidth}}
\toprule
\textbf{Preguntas / palabras clave} & \textbf{Notas / respuestas} \\
\midrule
\endfirsthead
\toprule
\textbf{Preguntas / palabras clave} & \textbf{Notas / respuestas} \\
\midrule
\endhead

\textbf{¿Qué es la anticresis y cómo se diferencia de la hipoteca?} &
Es el derecho real de garantía sobre cosas registrables (muebles o inmuebles). A diferencia de la hipoteca, hay desplazamiento de posesión al acreedor (o a un tercero designado por las partes). La publicidad la otorga la posesión y no el registro. \\[0.6em]

\textbf{¿Quiénes pueden constituir anticresis?} &
El titular del dominio, los condóminos, los propietarios en propiedad horizontal, el superficiario y el usufructuario (art.\ 2213). El usufructuario está limitado temporalmente por su propio derecho: si es vitalicio, hasta su muerte; si es a término, hasta el cumplimiento del plazo. \\[0.6em]

\textbf{¿Cuáles son los plazos máximos de la anticresis?} &
5 años para muebles registrables y 10 años para inmuebles, desde la constitución (art.\ 2218). No se prorrogaron de 20 a 35 años, como ocurrió con la inscripción hipotecaria. \\[0.6em]

\textbf{¿Qué derechos tiene el acreedor anticresista?} &
Usar y explotar la cosa, inclusive darla en arrendamiento, y percibir los frutos. En un inmueble, los frutos son civiles (rentas, alquileres) y se imputan primero a gastos e intereses y luego a capital. \\[0.6em]

\textbf{¿Qué deberes tiene el acreedor anticresista?} &
Conservar y mantener la cosa, afrontar los gastos de conservación y restituirla al constituyente al cumplirse el plazo o antes si se extingue la obligación. Puede recuperar el costo de las mejoras necesarias y el de las útiles solo si dieron a la cosa mayor valor. \\

\midrule
\multicolumn{2}{p{\dimexpr0.93\textwidth+2\tabcolsep\relax}}{%
\textbf{Resumen:} la anticresis recae sobre cosas registrables (muebles o inmuebles) e implica el desplazamiento posesorio al acreedor, quien puede usar la cosa y percibir sus frutos, con plazos máximos de 5 y 10 años según el bien. El acreedor debe conservar la cosa y restituirla.} \\
\bottomrule
\end{longtable}

% ============================================================
\chapter[Prenda]{Prenda}
% ============================================================

\section{Concepto legal}

\begin{definition}{Prenda (art.\ 2219 CCCN)}
``La prenda es el derecho real de garantía sobre cosas muebles no registrables o créditos instrumentados.''
\end{definition}

Debe distinguirse la \textbf{prenda clásica} (la regulada por el CCCN) de la \textbf{prenda con registro} (regulada por leyes especiales), que son dos figuras completamente distintas.

\section{Prenda clásica y prenda con registro}

\begin{table}[H]
\centering
\begin{tabularx}{\textwidth}{
    >{\raggedright\arraybackslash}p{0.17\textwidth}
    >{\raggedright\arraybackslash}X
    >{\raggedright\arraybackslash}X
}
\toprule
\textbf{Característica} & \textbf{Prenda clásica (CCCN)} & \textbf{Prenda con registro (leyes especiales)} \\
\midrule
Objeto & Cosas muebles \textbf{no} registrables o créditos instrumentados. & Cosas muebles registrables (autos, motos, maquinaria industrial, etc.). \\
Posesión & \textbf{Siempre} se desplaza al acreedor. & El constituyente conserva la posesión y puede seguir usando el bien. \\
Publicidad & La posesión del acreedor. & La inscripción en el registro prendario o en el registro del automotor. \\
Normativa aplicable & Arts.\ 2219 y ss.\ del CCCN. & Leyes especiales (Decreto-Ley 15.348/46 y normas del RNPA para automotores). \\
\bottomrule
\end{tabularx}
\end{table}

\begin{example}{Prenda con registro: automotor}
Cuando se compra un automotor y se paga en cuotas, el auto está prendado. La posesión la tiene el comprador, no la concesionaria que otorgó el préstamo ni el banco. El comprador es constituyente y a su vez tiene la posesión. La prenda registrable otorga publicidad por el registro, no por la posesión.
\end{example}

\begin{example}{Prenda con registro: industria}
Se utiliza mucho en la industria metalúrgica, donde lo que se prendan son máquinas. Esta prenda se realiza mediante su inscripción en un registro prendario. Lo importante es que el constituyente puede seguir utilizando las máquinas; de lo contrario, para obtener el crédito, el propietario debería entregar la posesión y carecería de ellas para producir.
% TODO: contenido incompleto o ambiguo en las notas originales (el apunte dice «el acreedor puede seguir utilizando las máquinas», pero el sentido del ejemplo indica que se refiere al constituyente)
\end{example}

\section{Legitimados y forma de constitución}

La prenda puede ser constituida por el dueño de la cosa.

\begin{important}
El contrato de prenda puede ser público o privado. En todos los casos es necesaria la \textbf{fecha cierta}. Para que se perfeccione la garantía es requisito indispensable (constitutivo) la \textbf{entrega de la cosa}, ya que el acreedor ejerce el derecho real por la posesión.
\end{important}

\begin{note}
A diferencia de lo que ocurre en la hipoteca, el condómino \textbf{no} puede prendar la cosa por su parte indivisa, porque tiene derechos muy restringidos respecto de toda la cosa y no podría entregarla. En la prenda, debe ser constituida necesariamente por todos los condóminos en conjunto.
\end{note}

\section{Prenda de créditos instrumentados}

\begin{definition}{Prenda de créditos instrumentados (art.\ 2232 CCCN)}
Existe la posibilidad de prender cualquier crédito instrumentado que pueda ser cedido. El caso típico es el pagaré.
\end{definition}

\begin{example}{Funcionamiento de la prenda de un pagaré}
Quien es acreedor de un pagaré (alguien le debe dinero y está documentado en ese instrumento) necesita un crédito. En lugar de ceder el pagaré, lo entrega en prenda a su acreedor como garantía. Para que quede constituida la garantía no basta la mera entrega del instrumento: se requieren \textbf{dos elementos constitutivos}:
\begin{enumerate}
    \item La entrega de la posesión del instrumento al acreedor prendario.
    \item La notificación al deudor cedido (quien firmó el pagaré), informándole que ya no debe pagar al constituyente sino directamente al acreedor prendario.
\end{enumerate}
\end{example}

El acreedor prendario, a partir de la constitución de la prenda, puede realizar todas las gestiones necesarias para cobrar el crédito: intimar el pago, ejecutar el pagaré, etc. Una vez extinguida la obligación garantizada, debe restituir el instrumento y rendir cuentas al constituyente de lo que percibió, para determinar si existe algún remanente que le corresponda a este último.

\section{Ejecución de la prenda y prenda de dinero}

\begin{definition}{Ejecución de la prenda (arts.\ 2228 y 2229 CCCN)}
Si el deudor no cumple, se realizará el bien en subasta pública conforme a lo regulado en los artículos 2228 y 2229 del Código Civil y Comercial.
\end{definition}

Se plantea el interrogante de si puede prendarse el dinero y si la moneda extranjera puede ser objeto de prenda, teniendo en cuenta que jurídicamente es una cosa mueble.

\begin{example}{Prenda de moneda extranjera en contratos de locación}
Es muy habitual que, cuando se celebra un contrato de locación y el locatario no tiene garantía propietaria, se constituya una prenda de moneda extranjera como garantía. La moneda extranjera, jurídicamente cosa mueble, queda prendada y es restituida al locatario a la finalización del contrato.
\end{example}

\subsection*{Cuadro Cornell: prenda}

\begin{longtable}{p{0.24\textwidth}p{0.69\textwidth}}
\toprule
\textbf{Preguntas / palabras clave} & \textbf{Notas / respuestas} \\
\midrule
\endfirsthead
\toprule
\textbf{Preguntas / palabras clave} & \textbf{Notas / respuestas} \\
\midrule
\endhead

\textbf{¿Qué es la prenda clásica y cuál es su objeto?} &
Es la prenda regulada por el CCCN (art.\ 2219). Recae sobre cosas muebles no registrables o créditos instrumentados. Siempre hay desplazamiento de posesión al acreedor. No debe confundirse con la prenda con registro (leyes especiales), que permite al constituyente conservar la posesión. \\[0.6em]

\textbf{¿Qué diferencia a la prenda clásica de la prenda con registro?} &
En la clásica, la posesión siempre se desplaza al acreedor y la publicidad la da la posesión. En la prenda con registro (automotores, maquinaria industrial), el constituyente conserva la posesión y la publicidad la otorga el registro. Las normas aplicables también difieren: el CCCN rige la primera y leyes especiales regulan la segunda. \\[0.6em]

\textbf{¿Quién puede constituir prenda y con qué forma?} &
El dueño de la cosa. El contrato puede ser público o privado, siempre con fecha cierta, y es requisito constitutivo la entrega de la cosa. Los condóminos no pueden prendar su parte indivisa: la prenda debe ser constituida por todos en conjunto. \\[0.6em]

\textbf{¿Cómo funciona la prenda de créditos instrumentados?} &
Se prendan créditos que constan en instrumentos cedibles (p.\,ej., pagarés). Se requieren dos elementos: (1) entrega de la posesión del instrumento al acreedor prendario; (2) notificación al deudor cedido, indicándole que el pago debe hacerlo al acreedor prendario. El acreedor puede cobrar el crédito y debe restituir el instrumento y rendir cuentas al constituyente. \\[0.6em]

\textbf{¿Cómo se ejecuta la prenda?} &
Si el deudor no cumple, el bien se realiza en subasta pública conforme a los arts.\ 2228 y 2229 CCCN. \\[0.6em]

\textbf{¿Puede prendarse el dinero o la moneda extranjera?} &
Sí. Jurídicamente la moneda extranjera es una cosa mueble. En la práctica se utiliza con frecuencia en contratos de locación cuando el locatario carece de garantía propietaria: se prenda moneda extranjera que se restituye al finalizar el contrato. \\

\midrule
\multicolumn{2}{p{\dimexpr0.93\textwidth+2\tabcolsep\relax}}{%
\textbf{Resumen:} la prenda clásica del Código Civil y Comercial recae sobre muebles no registrables o créditos instrumentados, exige siempre la entrega de la posesión al acreedor y no debe confundirse con la prenda con registro, figura especial regulada por leyes separadas en la que el constituyente conserva la posesión.} \\
\bottomrule
\end{longtable}

\end{document}
